% ============================================================
%  第三章 系统需求分析
% ============================================================
\chapter{系统需求分析}

\section{系统总体需求}

\subsection{用户与权限主体}

系统中“账号角色”与“社团内岗位”是两层概念。账号角色用于全局访问控制，决定用户在平台范围内的可访问资源；社团内岗位用于社团内部管理授权，决定成员在具体社团场景下的操作边界。未登录访客仅能浏览公开信息，不参与系统授权判断。为准确反映系统的双层权限结构，本文分别从系统角色与社团内部岗位两个层次进行说明，具体如表~\ref{tab:user-roles} 与表~\ref{tab:club-posts} 所示。

\begin{table}[H]
\centering
\caption{系统角色与权限范围}
\label{tab:user-roles}
\begin{tabular}{p{2.8cm}p{3.2cm}p{7.6cm}}
\toprule
系统角色 & 角色代码 & 权限范围 \\
\midrule
普通用户 & USER & 个人资料维护、社团浏览、活动报名与签到、提交招新申请、查看个人相关信息 \\
社团管理员 & CLUB\_ADMIN & 管理所属社团的活动、成员、招新、公告、资源申请、财务申请及相关社团信息维护 \\
系统管理员 & ADMIN & 执行全局审批、系统审计、统计分析、用户状态管理及平台级运维控制 \\
\bottomrule
\end{tabular}
\end{table}

\begin{table}[H]
\centering
\caption{社团内部岗位与操作范围}
\label{tab:club-posts}
\begin{tabular}{p{2.8cm}p{3.2cm}p{7.6cm}}
\toprule
社团岗位 & 岗位代码 & 操作范围 \\
\midrule
社长 & PRESIDENT & 负责社团整体组织与运行，具备社团内部管理权限，可处理成员、活动、公告、资源、财务和招新等核心事务 \\
管理员 & MANAGER & 协助社长完成社团日常事务管理，具备社团内部管理权限，可执行多数社团日常管理操作 \\
普通成员 & MEMBER & 主要具备信息查看、活动参与、报名申请与个人记录查询等基础操作权限 \\
\bottomrule
\end{tabular}
\end{table}

从实现机制上看，系统并未采用独立权限点表对每一项社团操作进行按钮级授权，而是采用“岗位角色判断 + 业务规则校验”的方式实现社团内部权限控制。当前实现中，\texttt{PRESIDENT} 与 \texttt{MANAGER} 在大多数业务场景下共享社团管理权限，统一纳入社团管理员能力范围；\texttt{MEMBER} 主要承担参与、查看和申请类操作。与此同时，系统还会结合社团状态、成员身份和具体业务流程对操作进行进一步限制，从而保证社团管理过程的规范性与安全性。
\subsection{系统边界}

本系统覆盖“社团创建与审批、成员管理、活动管理、招新管理、公告与通知、资源申请、财务审批、审计与统计、AI 问答与推荐”等完整流程。系统暂不包含真实支付结算、统一身份平台（CAS/OAuth）对接、移动端原生 App。

\subsection{需求来源与分析方法}

本文的需求分析主要服务于系统设计与实现型论文的工程建模，需求来源并非单一问卷或访谈结果，而是基于以下三类依据进行综合归纳：

\begin{enumerate}
  \item \textbf{典型业务流程}：围绕高校社团管理中较常见的创建审批、活动组织、招新审核、资源申请、财务审批与通知触达等流程，梳理系统需要覆盖的核心环节；
  \item \textbf{角色与权限边界}：结合访客、普通用户、社团管理员与系统管理员等主体的职责差异，归纳不同角色在认证、授权与业务操作上的功能诉求；
  \item \textbf{现有系统不足与项目约束}：结合绪论中的现有系统对比分析，以及本项目当前的部署条件、技术栈与实现范围，提炼非功能性需求与 AI 服务边界。
\end{enumerate}

因此，第三章更强调“面向实现的场景化需求整理”，其作用在于为后续系统设计、模块划分与测试验证提供依据，而不是将其作为独立的人因实验或问卷研究结论。

\section{功能性需求}

\subsection{认证与账户需求}

\begin{enumerate}
  \item 用户可注册、登录、登出、刷新会话、找回密码；
  \item 密码强度要求为“至少 8 位且同时包含字母和数字”；
  \item 登录失败计数基于 Redis，15 分钟窗口内超过 10 次将触发限频；
  \item 会话以 HttpOnly Cookie 传递，服务端可主动失效会话。
\end{enumerate}

\subsection{社团管理需求}

\begin{enumerate}
  \item 用户可提交社团创建申请，初始状态为 \texttt{PENDING}；
  \item 管理员审批通过后社团状态变为 \texttt{ACTIVE}；
  \item 社团支持解散流程：\texttt{DISSOLVING}（冷静期）到 \texttt{DISSOLVED}；
  \item 社团管理员可进行成员增删、角色调整、社团信息维护、成员导出。
\end{enumerate}

\subsection{活动管理需求}

\begin{enumerate}
  \item 社团管理员可创建活动并发布；
  \item 用户可报名活动，重复报名应被拒绝；
  \item 签到需满足“已报名 + 活动时间窗口有效 + 签到码正确（若配置）”；
  \item 支持签到明细导出与 WebSocket 实时签到人数广播；
  \item 活动状态支持 \texttt{DRAFT/PUBLISHED/IN\_PROGRESS/ENDED} 的自动推进。
\end{enumerate}

\subsection{招新管理需求}

\begin{enumerate}
  \item 社团管理员可创建招新批次与动态表单字段；
  \item 用户可提交申请并查询本人申请历史；
  \item 招新审核分一审、终审两阶段；
  \item 终审通过前需校验名额，且录取后自动补充成员关系；
  \item 支持批次申请数据导出（Excel）。
\end{enumerate}

\subsection{公告与通知、资源与财务需求}

\begin{enumerate}
  \item 支持站内通知分页查询、未读统计、单条已读与全部已读；
  \item 社团或系统公告发布后应触发异步通知分发并落库，前端可轮询读取最新通知状态；
  \item 资源申请需进行时段冲突/库存校验，并由管理员审批；
  \item 财务申请采用“提交-审批”流程，支出申请需校验社团余额。
\end{enumerate}

\subsection{AI 服务需求}

\begin{enumerate}
  \item 用户可通过聊天入口进行社团知识问答；
  \item 系统可按用户兴趣与行为推荐社团；
  \item 问答服务需要具备会话上下文能力与降级可用性；
  \item 推荐服务需具备冷启动兜底策略。
\end{enumerate}

\section{非功能性需求}

\begin{enumerate}
  \item \textbf{安全性}：接口访问遵循认证与授权规则，敏感操作可审计追踪；
  \item \textbf{一致性}：高并发写路径通过数据库锁控制，避免重复审批与重复报名；
  \item \textbf{可用性}：通知与部分非关键流程采用消息异步分发，减少主流程阻塞；
  \item \textbf{可维护性}：模块化单体组织结构，统一异常处理与统一响应格式；
  \item \textbf{可观测性}：暴露业务与运行指标，支持 Prometheus 抓取与 Grafana 展示。
\end{enumerate}

\paragraph{验证口径说明}
第三章提出的是面向实现的目标需求，而不是默认都已经以同等强度完成实证验证。后文对这些需求采用三类口径加以区分：一是已有直接测试证据支持的条目，如认证边界、并发写路径与部分性能基线；二是已有实现与局部观察支持、但尚缺专项实验的条目，如异步通知、可观测性与 AI 问答降级；三是仅完成实现基础、仍待更严格评估的条目，如推荐相关性质量与 AI 模块的外部可推广性。因此，后续章节不会将“已实现”直接表述为“已充分验证”，而是结合证据强度明确给出结论边界。

\section{用例分析}

\subsection{系统总用例图}

\begin{figure}[H]
\centering
\begin{tikzpicture}[
  actor/.style={draw, circle, minimum size=0.6cm},
  usecase/.style={draw, ellipse, minimum width=3cm, minimum height=0.8cm, align=center, font=\small},
  line/.style={-}
]
\draw[thick] (0,0) rectangle (13.6,9.4);
\node[above] at (6.8,9.4) {\textbf{校园社团管理系统}};

\node[actor] (visitor) at (-1.8,7.8) {};
\node[below] at (-1.8,7.4) {\small 访客};
\node[actor] (user) at (-1.8,4.9) {};
\node[below] at (-1.8,4.5) {\small USER};
\node[actor] (clubadmin) at (15.2,4.9) {};
\node[below] at (15.2,4.5) {\small CLUB\_ADMIN};
\node[actor] (admin) at (15.2,7.6) {};
\node[below] at (15.2,7.2) {\small ADMIN};

\node[usecase] (browse) at (3.2,7.8) {浏览社团/活动/公告};
\node[usecase] (auth) at (3.2,6.1) {登录/注册/找回密码};
\node[usecase] (apply) at (3.2,4.3) {活动报名/签到\\提交招新申请};
\node[usecase] (sysops) at (10.1,7.8) {全局审批与审计\\统计看板};
\node[usecase] (clubops) at (10.1,6.1) {社团管理\\成员与公告管理};
\node[usecase] (activityops) at (10.1,4.5) {活动发布\\招新审核};
\node[usecase] (resourceops) at (10.1,2.9) {资源与财务申请};
\node[usecase] (aiops) at (6.8,1.3) {AI 问答与推荐};

\draw[line] (visitor) -- (browse.west);
\draw[line] (visitor) -- (auth.north west);
\draw[line] (user) -- (auth.south west);
\draw[line] (user) -- (apply.west);
\draw[line] (user) -- ++(0.8,0) |- (aiops.west);

\draw[line] (clubadmin) -- (activityops.east);
\draw[line] (clubadmin) -- ++(-0.8,0) |- (clubops.east);
\draw[line] (clubadmin) -- ++(-1.2,0) |- (resourceops.east);
\draw[line] (clubadmin) -- ++(-1.6,0) |- (aiops.east);

\draw[line] (admin) -- (sysops.east);
\draw[line] (admin) -- ++(-4.0,0) |- (clubops.north east);
\draw[line] (admin) -- ++(-4.4,0) |- (resourceops.north east);
\end{tikzpicture}
\caption{系统总用例图}
\label{fig:usecase}
\end{figure}

\subsection{核心用例：活动签到}

\begin{table}[H]
\centering
\caption{活动签到用例描述}
\begin{tabular}{lp{9.2cm}}
\toprule
项目 & 描述 \\
\midrule
用例名称 & 活动签到 \\
参与者 & 已登录用户（通常为已加入社团成员） \\
前置条件 & 已报名活动，当前时间位于活动开始与结束时间之间 \\
主成功场景 &
1) 用户提交签到请求并携带签到码（若活动设置） \newline
2) 系统校验报名状态与时间窗口 \newline
3) 系统校验签到码并检测是否重复签到 \newline
4) 写入签到记录并更新报名状态 \newline
5) 通过 WebSocket 广播最新签到人数 \\
异常分支 &
2a) 未报名：拒绝签到 \newline
2b) 活动未开始或已结束：拒绝签到 \newline
3a) 签到码错误：拒绝签到 \newline
3b) 已签到：提示重复签到 \\
后置条件 & 签到数据入库，前端看板可实时看到人数变化 \\
\bottomrule
\end{tabular}
\end{table}

\begin{figure}[H]
\centering
\begin{tikzpicture}[
  node distance=0.9cm,
  step/.style={draw, rectangle, rounded corners, minimum width=4.2cm, minimum height=0.78cm, align=center, font=\small},
  decision/.style={draw, diamond, aspect=2.3, inner sep=1pt, align=center, font=\small},
  line/.style={-Stealth, thick}
]
\node[step] (start) {用户提交签到请求};
\node[decision, below=of start] (timecheck) {是否已报名且处于\\签到时间窗口?};
\node[decision, below=of timecheck] (codecheck) {签到码是否正确\\或无需签到码?};
\node[decision, below=of codecheck] (duplicate) {是否已经签到?};
\node[step, below=of duplicate] (save) {写入签到记录并更新报名状态};
\node[step, below=of save] (push) {广播最新签到人数并返回成功};

\node[step, right=4.2cm of timecheck] (rejecttime) {拒绝签到\\返回未报名/时间不合法原因};
\node[step, right=4.2cm of codecheck] (rejectcode) {拒绝签到\\返回签到码错误};
\node[step, right=4.2cm of duplicate] (already) {返回已签到结果\\不重复写入数据};

\draw[line] (start) -- (timecheck);
\draw[line] (timecheck) -- node[left,font=\tiny]{是} (codecheck);
\draw[line] (timecheck.east) -- ++(0.7,0) |- node[pos=0.25,above,font=\tiny]{否} (rejecttime.west);
\draw[line] (codecheck) -- node[left,font=\tiny]{是} (duplicate);
\draw[line] (codecheck.east) -- ++(0.7,0) |- node[pos=0.25,above,font=\tiny]{否} (rejectcode.west);
\draw[line] (duplicate) -- node[left,font=\tiny]{否} (save);
\draw[line] (duplicate.east) -- ++(0.7,0) |- node[pos=0.25,above,font=\tiny]{是} (already.west);
\draw[line] (save) -- (push);
\end{tikzpicture}
\caption{活动签到业务流程图}
\label{fig:activity-checkin-flow}
\end{figure}

图~\ref{fig:activity-checkin-flow} 对应了表格用例中的主成功场景与异常分支：系统先校验报名状态和时间窗口，再校验签到码，最后在防重条件下完成签到记录写入与 WebSocket 广播。这一流程图有助于将“业务规则约束”与后续设计章节中的接口时序、并发控制和实时推送设计对应起来。
